\section{A fixed Gaussian does not give a bounded ideal inverse}
\label{sec:dense-gaussian-inverse-obstruction}

\paragraph{Status and scope.}
The complete post-writing author reading and two independent
whole-proof readings have passed.
MA-PI-013 constructs a preimage for each smooth theta
vector.  The following shows that these preimages cannot
satisfy a common estimate in the exact central printed
Hilbert norm.  The obstruction occurs at one established
critical-line zero, even at fixed Gaussian scale.

Use the fixed seed, \(R\), \(B\), \(\Psi\), and \(\psi\)
of MA-PI-013, and set
\[
 \begin{gathered}
 H_s=H^{\rm one}_{0,s}
       =L^2(\mathbb R,(1+x^2)^{s/2}\,dx),\qquad s\ge0,\\
 J_s=\overline{\mathscr I}^{\,H_s},\qquad
 \gamma_\tau(x)=\frac1{\sqrt{4\pi\tau}}
                       e^{-x^2/(4\tau)},\qquad \tau>0,\\
 \mathcal M\gamma_\tau(z)=e^{\tau z^2}.
 \end{gathered}
 \tag{MA-PI.320}
\]
This is the actual printed \(A=1/2\) realization, after
the exact Haar/unitary identification of MA-PI-010.
It is not a positive-width symmetric norm.

\paragraph{MA-PI-014 (failure of a uniform Gaussian preimage estimate).}
For every \(s\ge0\) and every fixed \(\tau>0\), there is
no bounded linear map \(S_\tau:J_s\to H_s\) such that
\[
                  RS_\tau a=C_{\gamma_\tau}a
                       \qquad(a\in J_s).
 \tag{MA-PI.321}
\]
Indeed there is no constant \(C\) giving the corresponding
bound even on the actual smooth ideal:
\[
 \|S_\tau a\|_{H_s}\le C\|a\|_{H_s},
       \qquad a\in\mathscr I,\quad
       RS_\tau a=C_{\gamma_\tau}a.
 \tag{MA-PI.322}
\]
The smooth-ideal preimages here exist by MA-PI-013:
replace \(e^{z^2/n^2}\) by \(e^{\tau z^2}\) in the
same Gaussian formula and contour proof.
They are unique because \(R\) is injective.  The theorem
therefore concerns their actual norms, not a failure
to choose a convenient linear section.

\begin{proof}
\emph{1. The known zero and its exact local factor.}
Choose a nontrivial critical-line zero
\(\rho=1/2+iT\) of zeta, let \(\lambda=iT\), and
let \(m\ge1\) be its finite multiplicity.  Such zeros
are known unconditionally.  For example, the
published input recorded in MA-MO-023, Conrey's
Theorem 1 with its counting conventions, supplies
infinitely many simple choices.  No assertion about
the remaining zeros is made.

The entire function \(\Psi\) has the same multiplicity
at \(\lambda\), since its other factors are nonzero
there.  Put
\[
       \Phi_\lambda(z)=\frac{\Psi(z)}{(z-\lambda)^m}.
 \tag{MA-PI.323}
\]
The apparent singularity is removable.
On every fixed vertical strip, this entire function
and its fixed derivatives decrease exponentially in
height times a polynomial, just as \(\Psi\) does.
Outside a fixed disk around \(\lambda\) this follows
by division by the polynomial; on that disk it
follows from holomorphy and compactness.
The complete inverse-Mellin argument in MT-M18.6
therefore gives a kernel \(\varphi_\lambda\in\mathcal E\)
with transform \(\Phi_\lambda\).
For every compact smooth \(w\),
\[
 C_\psi w
   =C_{\varphi_\lambda}(-D-\lambda)^m w,
                       \qquad D=d/dx.
 \tag{MA-PI.324}
\]
Indeed the Mellin multiplier of \(-D-\lambda\)
is \(z-\lambda\).  Fourier uniqueness proves
the actual convolution identity.

\emph{2. Compact packets and the precise polynomial weight.}
Fix nonzero \(\chi\in C_c^\infty((1,2))\) and,
for \(0<\epsilon\le1\), set
\[
 w_\epsilon(x)=e^{-\lambda x}\chi(\epsilon x),
       \qquad y_\epsilon=C_\psi w_\epsilon\in\mathscr I.
 \tag{MA-PI.325}
\]
The inputs \(w_\epsilon\) are genuinely compact
smooth functions.  Since \(\lambda\) is imaginary,
\[
 \|w_\epsilon\|_{H_s}^2
   =\epsilon^{-s-1}
       \int_1^2|\chi(t)|^2(\epsilon^2+t^2)^{s/2}\,dt
   \asymp_{s,\chi}\epsilon^{-s-1}.
 \tag{MA-PI.326}
\]
Replacing \(\chi\) by any fixed derivative gives
the corresponding upper bound.  Thus
\[
 \|e^{-\lambda\cdot}\chi^{(j)}(\epsilon\cdot)\|_{H_s}
             \le C_{s,\chi,j}\|w_\epsilon\|_{H_s}.
 \tag{MA-PI.327}
\]
All constants here are independent of \(\epsilon\).
Equation (MA-PI.324) and
\[
 (-D-\lambda)^m w_\epsilon
       =(-\epsilon)^m e^{-\lambda x}
                              \chi^{(m)}(\epsilon x)
 \tag{MA-PI.328}
\]
give
\[
          \|y_\epsilon\|_{H_s}
              \le C_{s,\chi,\lambda}\epsilon^m
                                      \|w_\epsilon\|_{H_s}.
 \tag{MA-PI.329}
\]
Here \(C_{\varphi_\lambda}\) is bounded on \(H_s\)
by the actual weighted Young estimate, since
\(\varphi_\lambda\in\mathcal E\).

The vectors \(y_\epsilon\) are actual original
source vectors, not merely functions satisfying
some vanishing conditions.  They belong to
\(\psi*\mathcal E\subset\mathscr I\).
More explicitly, MT-M18.4 and MT-M24.7--9,
with the one-place absorption of MA-TR-101,
give a fixed even Schwartz \(f_\psi\) with
\(f_\psi(0)=\int f_\psi=0\) and
\(\Theta f_\psi=\psi\).
Then the compact-parameter integral
\[
 f_\epsilon=\int_{\mathbb R}
                 w_\epsilon(t)D_t f_\psi\,dt
       \in\mathscr X_0,\qquad
       \Theta f_\epsilon=y_\epsilon
 \tag{MA-PI.330}
\]
uses exactly the normalized dilation of MA-PI.250.
The integral converges in every Schwartz seminorm
for each fixed \(\epsilon\); both polar conditions
are preserved termwise.  No uniform bound on these
source seminorms is assumed.

\emph{3. The forced preimages retain a nonzero leading multiplier.}
Let \(\ell\) be the actual kernel of \(B\), and put
\[
 \kappa_\tau=\ell*\gamma_\tau,\qquad
 c_\tau=\mathcal M\kappa_\tau(\lambda)
        =\mathfrak b(\lambda)e^{\tau\lambda^2}\ne0.
 \tag{MA-PI.331}
\]
Nonvanishing follows directly from (MA-PI.286):
\(\lambda+1/2\ne0\), both indicated gamma factors
are finite and nonzero, and \(H(\lambda)\ne0\).
In particular the Gaussian never cancels this
local value.  The kernel \(\kappa_\tau\) is
Schwartz with every polynomial moment integrable.

Set \(q_\epsilon(x)=\chi(\epsilon x)\).
The exact weight inequality and the fundamental
theorem of calculus give
\[
 \begin{aligned}
 \|\tau_y q_\epsilon-q_\epsilon\|_{H_s}
 &\le C_s|y|(1+y^2)^{s/4}\|q_\epsilon'\|_{H_s}\\
 &\le C_{s,\chi}\epsilon |y|(1+y^2)^{s/4}
                                      \|w_\epsilon\|_{H_s}.
 \end{aligned}
 \tag{MA-PI.332}
\]
For negative \(y\), integrate the derivative over
the interval between \(y\) and zero; the same
bound applies.  Multiplication by \(e^{-\lambda x}\)
is isometric on \(H_s\).  Consequently the identity
\[
 \begin{split}
 C_{\kappa_\tau}w_\epsilon-c_\tau w_\epsilon
   =e^{-\lambda x}\int \kappa_\tau(y)e^{\lambda y}
                [q_\epsilon(x-y)-q_\epsilon(x)]\,dy
 \end{split}
 \tag{MA-PI.333}
\]
and (MA-PI.332) imply
\[
 \|C_{\kappa_\tau}w_\epsilon-c_\tau w_\epsilon\|_{H_s}
       \le C_{\tau,s,\chi}\epsilon\|w_\epsilon\|_{H_s}.
 \tag{MA-PI.334}
\]
The integral of
\(|\kappa_\tau(y)||y|(1+y^2)^{s/4}\)
is finite.  Hence, for sufficiently small
\(\epsilon\), depending on the fixed parameters,
\[
       \|C_{\kappa_\tau}w_\epsilon\|_{H_s}
                    \ge \tfrac12|c_\tau|
                                      \|w_\epsilon\|_{H_s}.
 \tag{MA-PI.335}
\]

By the exact factorization \(RB=C_\psi\) and
commutativity of the bounded convolution maps,
\[
 R(C_{\kappa_\tau}w_\epsilon)
       =C_{\gamma_\tau}y_\epsilon .
 \tag{MA-PI.336}
\]
Injectivity of \(R\) forces every preimage in
(MA-PI.322) to equal this vector.  The numerator
in (MA-PI.335) is nonzero, and \(y_\epsilon\ne0\):
alternatively the latter follows directly from
the nonzero entire product
\(\Psi\,\mathcal Mw_\epsilon\).
Combining (MA-PI.329) and (MA-PI.335) gives
\[
 \frac{\|S_\tau y_\epsilon\|_{H_s}}
      {\|y_\epsilon\|_{H_s}}
       \ge c_{\tau,s,\chi,\lambda}\epsilon^{-m}
                         \longrightarrow\infty .
 \tag{MA-PI.337}
\]
Normalizing the genuine smooth-ideal inputs proves
both assertions.  In particular
\(\tau=1/n^2\) rules out a bounded extension
of each fixed PI013 Gaussian preimage map.
\end{proof}

\paragraph{Exact implication for the proposed route.}
The failure is a local division loss at one fixed
zero, not an estimate about reciprocal growth at
increasing height.  No polynomial bound in the
zero ordinate could remove (MA-PI.337).
It does not contradict MA-MO-028: that return
contains the actual local commutator
\(C_g\mathcal L\), uses a positive-width
two-sided norm, and proves a different operator
identity on the completed ideal.  It also does
not disprove the full-source criterion MA-MO-036
or its conditional sufficient cases.
No positive-width symmetric or endpoint extension
of the present obstruction is asserted.

The known-zero input is the published partial
zero-location theorem already documented in
MA-MO-023: Conrey (1989), Theorem 1, printed
p.~4, with equations (8), (11), and (12),
printed p.~3.  Only existence of one
critical-line zero is needed here; no fresh
audit of the mollifier proof is claimed.

Conrey, J. B. (1989). More than two fifths of
the zeros of the Riemann zeta function are on
the critical line. \emph{Journal f\"ur die
reine und angewandte Mathematik, 399}, 1--26.
\url{https://doi.org/10.1515/crll.1989.399.1}.

The local multiplicity \(m\) is the length of
the corresponding finite divisor factor.
This algebraic factor gives the exact
\(\epsilon^m\) norm loss, while Gaussian
smoothing is a unit at that factor.  Neither
finite-length bookkeeping nor dense range
repairs the analytic norm failure.  This
decides the proposed Gaussian estimate in
the stated metric, not a previously open
FT002 arithmetic gate.
